\chapter{First Aid and CPR}\label{ch:first-aid-cpr}

The candidate must know what first aid is, the four steps in any emergency, how to check breathing, and the main figures of adult CPR. This chapter is an introduction for the exam. It does not replace a hands-on first-aid course.

\begin{caution}
Medical content. A qualified reviewer must check this chapter before print. The emergency number 112 is still to be confirmed on an official page. The figures follow the Indian First Aid Manual, seventh edition (2016); re-read the current edition before print.
\end{caution}

\section*{What first aid is}

First aid is the first help given to an injured or suddenly ill person, before an ambulance, a medical person or a healthcare facility can take over.

Its aims:
\begin{itemize}
\item to preserve life;
\item to prevent the condition from getting worse;
\item to promote recovery;
\item to help the person reach the nearest healthcare facility safely.
\end{itemize}

\section*{Four steps in any emergency}

\begin{enumerate}
\item Make the area safe. Your own safety comes first.
\item Evaluate the person. Is the person conscious? Is the person breathing normally?
\item Seek help. Send a bystander to call, and to come back and confirm.
\item Give first aid.
\end{enumerate}

When calling for help, have ready:
\begin{itemize}
\item where you are;
\item the number you are calling from;
\item what has happened;
\item how many people are hurt, and how.
\end{itemize}

\section*{Moving a person, food and drink}

\begin{itemize}
\item As a rule, do not move an injured person.
\item Move the person only if staying is more dangerous, the place cannot be made safe, and help will not come soon. Even then, move only if you are not in danger yourself.
\item Give nothing to eat or drink to a person who is badly hurt, feels sick, is getting sleepy or is losing consciousness.
\item The manual lists a few exceptions, for example low blood sugar, only when the person is conscious.
\end{itemize}

\section*{Checking breathing}

\begin{itemize}
\item Tilt the head back and lift the chin.
\item Look, listen and feel for breathing for no more than 10 seconds.
\item Do not spend time feeling for a pulse.
\item Start CPR when the person is not breathing, or not breathing normally.
\item Occasional noisy gasps just after a collapse are not normal breathing.
\end{itemize}

\section*{CPR for an adult}

\subsection*{Chest compressions}

\begin{itemize}
\item Place the person on their back on a hard surface.
\item Kneel beside the upper arm.
\item Put the heel of one hand in the centre of the chest, the other hand on top, fingers locked.
\item Keep the shoulders over the chest and the arms straight.
\item Push down 5 to 6 cm.
\item Release fully without leaning, so that the chest rises each time.
\item Rate: about 100 compressions a minute, and not more than 120.
\item Give 30 compressions in a set.
\end{itemize}

\subsection*{Rescue breaths}

\begin{itemize}
\item After 30 compressions, tilt the head, lift the chin, pinch the nose and seal the mouth.
\item Blow for one second and watch the chest rise.
\item Give two breaths, with at most two attempts, then go back to 30 compressions. The ratio is 30 to 2.
\end{itemize}

\subsection*{Who gives breaths}

\begin{itemize}
\item An untrained or little trained helper gives compressions only.
\item A trained helper gives compressions with breaths.
\item For a baby under one year, breaths are always given, with the technique for babies.
\end{itemize}

\subsection*{Swapping, stopping, afterwards}

\begin{itemize}
\item If more than one rescuer is present, swap every two minutes, with the smallest possible pause.
\item Stop CPR only when the person shows signs of life such as breathing, a trained person or medical staff takes over, you are too exhausted, or the place becomes unsafe.
\item If breathing returns, place the person in the recovery position.
\item Wash your hands afterwards.
\end{itemize}

\section*{Role of the yoga instructor}

Call for help early. Give first aid or CPR only as trained. Hand over to medical staff. Learn CPR in a hands-on course of the Indian Red Cross or St John Ambulance.

The first aid kit and the steps for wounds, burns and fractures are not covered here. They are not yet checked against the manual.

\begin{recall}
\begin{itemize}
\item Order of the steps: safe, check, call, help.
\item Breathing check: up to 10 seconds. No pulse check.
\item Depth: 5 to 6 cm. Rate: 100 to 120 a minute. Ratio: 30 to 2.
\item Swap rescuers every two minutes.
\item Do not move an injured person unless staying is more dangerous.
\item Untrained helper: compressions only.
\end{itemize}
\end{recall}

\begin{caution}
Safety points. Your own safety first. Do not feel for a pulse. Do not give food or drink to a badly hurt or drowsy person. Do not move an injured person without a strong reason. Give CPR only as trained.
\end{caution}

\subsection*{Worked examples}

\textbf{The four steps in any emergency.} Give the steps in order and add one line on each: make the area safe, your own safety first; evaluate the person, conscious and breathing normally; seek help, sending a bystander to call and to come back and confirm; give first aid. Memory aid: safe, check, call, help.

\textbf{A person collapses in the hall.}
\begin{itemize}
\item Check that the place is safe for you and for the person.
\item Tilt the head back and lift the chin. Look, listen and feel for breathing for no more than 10 seconds. Do not spend time on a pulse.
\item Send a bystander to call for help and to come back and confirm. The caller has ready where you are, the number being called from, what has happened, and how many are hurt and how.
\item If the person is not breathing, or only gasps noisily just after the collapse, start CPR.
\item Push in the centre of the chest, 5 to 6 cm deep, about 100 a minute and not more than 120, in sets of 30. Add two breaths only if you are trained.
\item Stop only when there are signs of life such as breathing, a trained person or medical staff takes over, you are too exhausted, or the place becomes unsafe.
\end{itemize}

\textbf{Memory aid for the CPR numbers.} Breathing check: up to 10 seconds. Depth: 5 to 6 cm. Rate: 100 to 120 a minute. Ratio: 30 to 2. Swap rescuers every two minutes. Say the five numbers aloud in this order until they come without effort.

\subsection*{Common questions}

\begin{itemize}
\item \emph{Is the emergency number in India 112?} The notes give 112, but it is not in the manual and is still to be confirmed on an official page.
\item \emph{Will this chapter make me able to do CPR?} No. Learn it in a hands-on course.
\item \emph{Do the Indian manual and the American Heart Association give different figures?} No. The rate, depth and 30 to 2 ratio agree. Use the Indian manual for the exam.
\item \emph{Is CPR certification compulsory for yoga teachers in India?} The workshop deck says it is not, but this is still to be checked. A course is worth taking in any case.
\end{itemize}

\subsection*{Check yourself}

\begin{enumerate}
\item What is the emergency number in India, and what is still to be done about it in this book?
\item What compression rate does current American Heart Association guidance give for adult CPR?
\item What is the usual ratio of compressions to breaths for a trained helper?
\item Fill in the blank: the ratio of chest compressions to rescue breaths is \underline{\hspace{1cm}} to \underline{\hspace{1cm}}.
\item What is the compression rate for an adult?
\item True or false: you should feel for a pulse for a full minute before starting CPR.
\item Why is a noisy gasp just after a collapse not a sign to wait?
\item Who gives rescue breaths, and who gives compressions only?
\item A person has fallen and is badly hurt but is breathing, and the place is safe. May you move them to another spot just to make them comfortable?
\item A conscious person who feels sick asks for a drink after a fall. What do you do?
\item Two helpers are doing CPR and one is tiring. What should they do?
\end{enumerate}

\subsection*{Answers}

\begin{enumerate}
\item 112. It is not in the manual and is still to be confirmed on an official page.
\item 100 to 120 compressions a minute.
\item 30 compressions to 2 breaths.
\item 30 to 2.
\item About 100 a minute, and not more than 120.
\item False. Check breathing for no more than 10 seconds and do not spend time feeling for a pulse.
\item Occasional noisy gasps just after a collapse are not normal breathing, so CPR should start.
\item A trained helper gives compressions with breaths. An untrained or little trained helper gives compressions only. For a baby under one year, breaths are always given, with the technique for babies.
\item As a rule, no. Move only if staying is more dangerous, the place cannot be made safe and help will not come soon, and only if you are not in danger yourself.
\item Give nothing to eat or drink to a person who is badly hurt, feels sick, is getting sleepy or is losing consciousness. The manual lists a few exceptions, for example low blood sugar, only when the person is conscious.
\item Swap rescuers every two minutes with the smallest possible pause. If breathing returns, place the person in the recovery position and wash hands afterwards.
\end{enumerate}

\paragraph{Sources.} Indian Red Cross Society and St John Ambulance Association India, \emph{Indian First Aid Manual}, seventh edition (2016), the chapter on dealing with an emergency and basic CPR, cross-checked with American Heart Association guidance on CPR (2018 update); the Board syllabus for Level 1, Yoga Protocol Instructor; the official Ministry of AYUSH and Quality Council of India guidebook for Level I (older scheme, used as a coverage check). All in own words.
